\section{机器学习模型}

为系统评估不同架构在纳米毒性预测任务上的适用性，我们尝试了多种模型，涵盖传统树模型、梯度提升树、注意力增强树模型以及深度表格模型。经过初步性能比较，最终选择随机森林、XGBoost、ABRF、LARF、TabTransformer 以及蒸馏 TabTransformer 六种模型作为主要分析对象。下面将依次阐述各模型的设计思想及其在本任务中的定位。

随机森林是以决策树为基础的典型的集成学习方法，先以自助法抽样构造若干决策树再据此做回归或分类，因而天然具备抗过拟合能力强、能自然、合理地处理混合型特征（数值型和类别型）、对缺失值不敏感、能直接估计特征重要性多个优点，因此很自然地成为中小型表格数据的良好基线模型。故本文采用随机森林来系统、严谨地考察其在样本量有限的硒纳米颗粒数据集上的稳健性及可解释性。

XGBoost 是一种梯度提升树框架，把二阶导数信息及正则化项（L1、L2）都自然、合理地用上了，因此在各类表格数据竞赛中都有极好的表现。XGBoost 以串行方式迭代拟合残差，故一般能获得比随机森林更高的预测精度，但也正因如此，使用 XGBoost 时必须妥善调整超参数以避免过拟合。因此本文引入 XGBoost 的目的很明确：考察梯度提升树在引入增强特征（IC50、元素组成等）之后能否达到更优的预测性能。

ABRF 是在标准随机森林基础上加入了一个注意力网络来对各棵决策树的预测做动态加权处理：先训练普通的随机森林，再对每个样本收集所有树的预测值，接着用一个小型神经网络直接、自动地学习各棵树就该样本的权重，最后用各树预测值的加权和作为最终预测结果。因此其本质是为了用注意力机制合理地挖掘不同树、不同样本之间预测可靠性的差异。故本文引入 ABRF 的目的是考察单层注意力对树模型的实际增益。

LARF 在 ABRF 的基础上引入了特征级注意力：即在每棵决策树构造以前先学习一个特征权重向量对原始特征做加权处理，因此树分裂时更有利于利用重要特征，继而叠加树级注意力网络，由此自然地形成双层注意力机制，也因而能直接、明确地优化特征选择及树集成两个环节。故而本文引入 LARF 来考察更复杂的注意力结构是否有利于提高预测性能。

TabTransformer 是一种针对表格数据的深度学习方法，其核心思路是将类别特征通过 Transformer 的注意力机制编码为嵌入向量，再与数值特征拼接后输入多层感知机。Transformer 擅长捕捉类别特征之间的高阶交互关系，因此在包含多个分类变量的数据集上具有潜在优势。本研究将 TabTransformer 作为深度表格模型的代表，用于评估深度学习在纳米毒性数据上的性能，并与传统树模型进行对比。

蒸馏 TabTransformer 是用知识蒸馏方法将完整的 TabTransformer 作为教师模型，自然、合理地训练参数量更小的学生模型（即降低嵌入维数及 MLP 隐藏层尺寸），且所设计的蒸馏损失把教师输出的软标签及真实硬标签都加以利用，因此其目的是在保证较好精度的基础上得到轻量化模型。故本文引入该模型，就是为了系统、严谨地考察资源受限场景下的部署问题，即使性能略有下降，也具有极好的参考价值。

由于在大规模多类型纳米颗粒上对所列六种模型做了十分完备、有层次的比较，因此本文考察了不同架构对异质纳米颗粒数据的适应性。具体而言，在硒纳米颗粒专用数据集上，因样本量很小（383 条），故重点比较随机森林、XGBoost 和 TabTransformer，并系统分析基础特征及增强特征两种设置下的性能差异，而 ABRF、LARF 和蒸馏 TabTransformer 仅在基准数据集上作了评价，以兼顾计算成本及分析简洁性。

\section{评价指标说明}

本文对模型性能采用回归、分类两类指标进行评价：回归任务中采用决定系数（\(R^2\)）、平均绝对误差（MAE）和均方根误差（RMSE）。具体而言，\(R^2\) 是对模型所拟合的目标变量方差的解释比例的自然度量，取值范围为 0--1，越接近 1 拟合效果越好。MAE 是预测值与真实值绝对误差的算术平均，其单位与目标变量一致（\%），因此也是预测典型偏差的直接度量。RMSE 对较大误差更敏感，很自然地适用于考察预测的稳定性。

分类任务以 85\% 细胞存活率为阈值，将样本合理划分为“高存活”和“低存活”两类，继而用准确率（Accuracy）、F1 分数和 ROC 曲线下面积（AUC）作为主要指标。准确率即正确分类的比例，F1 分数是精确率、召回率的调和平均，对类别不平衡问题有极好的适应性，而 AUC 直接、客观地衡量了模型区分正负类的能力，0.5 为随机，1.0 为完美。

\section{结果与讨论}

\subsection{多类型纳米颗粒基准数据集上的模型比较}
为获得可靠的性能基线，我们首先在大规模多类型纳米颗粒毒性数据集上比较六类模型的表现，包括随机森林、XGBoost、ABRF、LARF、TabTransformer 以及蒸馏 TabTransformer。该基准数据集来源于 Masarkar 等人建立的公开纳米颗粒毒性数据资源 \cite{masarkar2024}，覆盖金属、金属氧化物、聚合物、碳基与无机纳米颗粒，清洗后共包含 4102 个有效样本，按照 7:3 划分为训练集 3281 条与测试集 821 条。输入特征包括浓度、暴露时间、材料类别、粒径、Zeta 电位、细胞类型与检测方法等。

回归任务结果如表~\ref{tab:benchmark_regression} 所示。TabTransformer 取得最高的决定系数，$R^2$ 为 0.768，MAE 为 10.30\%，RMSE 为 14.97\%。XGBoost 与随机森林表现接近，分别达到 $R^2=0.752$ 与 $R^2=0.749$。引入注意力机制后的 ABRF 与 LARF 未超过传统树模型，蒸馏 TabTransformer 的性能则明显下降。

\begin{table}[htbp]
  \centering
  \caption{多类型纳米颗粒基准数据集上的回归性能比较}
  \label{tab:benchmark_regression}
  \begin{tabular}{lccc}
    \hline
    模型 & $R^2$ & MAE (\%) & RMSE (\%) \\
    \hline
    TabTransformer & 0.768 & 10.30 & 14.97 \\
    XGBoost & 0.752 & 11.13 & 15.56 \\
    随机森林 & 0.749 & 10.81 & 15.71 \\
    ABRF & 0.719 & 11.87 & 16.55 \\
    LARF & 0.716 & 11.94 & 16.63 \\
    蒸馏 TabTransformer & 0.506 & 16.85 & 22.31 \\
    \hline
  \end{tabular}
\end{table}

在硒纳米颗粒专用数据集上，我们采用更贴近实际应用的三折交叉验证方案，并以基础特征与增强特征两种设置分别训练随机森林、XGBoost 与 TabTransformer。基础特征包括浓度、暴露时间、粒径、Zeta 电位、尺寸比、表面电荷、细胞类型、修饰基团和检测方法等；增强特征在此基础上进一步引入 IC50、元素 Se 含量、长期稳定性、结合能 Se3d 以及晶型编码。

结果如表~\ref{tab:senp_main_results} 所示。基础特征下，随机森林与 XGBoost 已具备较好的回归能力，分别达到 $R^2=0.701$ 与 $R^2=0.751$；TabTransformer 在该中小样本场景下明显欠拟合，$R^2$ 仅为 0.003。引入增强特征后，树模型性能显著提升：随机森林的三折 CV MAE 从 11.04\% 降至 7.63\%，$R^2$ 从 0.701 提升至 0.915；XGBoost 的三折 CV MAE 从 9.87\% 降至 7.20\%，$R^2$ 进一步提升至 0.957，为所有方案中最佳。相比之下，TabTransformer 即使加入增强特征后也仅达到 $R^2=0.057$，说明在当前样本规模下，树模型对硒纳米颗粒表格数据更具优势。

\begin{table}[htbp]
  \centering
  \caption{硒纳米颗粒数据集上基础特征与增强特征的模型性能}
  \label{tab:senp_main_results}
  \begin{tabular}{lccc}
    \hline
    模型 & 特征集 & 3折 CV MAE (\%) & $R^2$ \\
    \hline
    随机森林 & 基础 & 11.04 & 0.701 \\
    随机森林 & 增强 & 7.63 & 0.915 \\
    XGBoost & 基础 & 9.87 & 0.751 \\
    XGBoost & 增强 & 7.20 & 0.957 \\
    TabTransformer & 基础 & 24.13 & 0.003 \\
    TabTransformer & 增强 & 25.50 & 0.057 \\
    \hline
  \end{tabular}
\end{table}

\subsection{分类模型结果}

分类任务方面，树模型表现更稳健，其中 XGBoost 的 AUC 最高，为 0.945；随机森林的准确率与 F1 值最好，分别为 84.9\% 与 0.840。该基准比较表明，在大样本异质表格数据上，深度表格模型和树模型都具有较强竞争力，但树模型在分类任务中更稳定、训练成本更低。

\begin{table}[htbp]
  \centering
  \caption{数据集连续型特征的描述性统计（各特征可用样本量见 $n$ 列）}
  \label{tab:statistics}
  \begin{tabular}{lcccc}
    \hline
    模型 & Accuracy & F1 & AUC \\
    \hline
    随机森林 & 0.849 & 0.840 & 0.943 \\
    XGBoost & 0.849 & 0.836 & 0.945 \\
    TabTransformer & 0.845 & 0.835 & 0.939 \\
    ABRF & 0.828 & 0.817 & 0.921 \\
    LARF & 0.819 & 0.809 & 0.914 \\
    蒸馏 TabTransformer & 0.722 & 0.707 & 0.801 \\
    \hline
  \end{tabular}
\end{table}

\subsubsection{特征效应}
增强特征模型的特征重要性分析表明，IC50、元素 Se 含量、长期稳定性和结合能 Se3d 是影响硒纳米颗粒细胞毒性预测结果的核心变量。其中，IC50 直接反映剂量--效应关系，元素 Se 含量与颗粒纯度及潜在离子释放背景相关，长期稳定性表征颗粒在分散体系中的聚集与氧化趋势，结合能 Se3d 则反映表面化学态差异。该结果说明，在单一材料体系的毒性建模中，能够刻画材料本征化学属性的增强特征，比单纯依赖粒径、浓度和暴露时间等通用参数更具解释力。

\subsection{组合特征分析}
为了识别真正有效的增强信息，我们对五个候选增强特征进行了自动化筛选，候选特征包括 IC50\_uM、元素 Se 含量（wt\%）、长期稳定性（周）、结合能 Se3d（eV）与晶型编码。穷举 31 种非空组合并以三折交叉验证 MAE 为目标后，最优子集为
\[
\{\mathrm{IC50\_uM},\ \mathrm{元素\ Se\ 含量},\ \mathrm{长期稳定性},\ \mathrm{结合能\ Se3d}\}.
\]
该组合的三折 CV MAE 为 7.60\%，与使用全部增强特征时的 7.63\% 基本一致，说明晶型编码在当前数据集中的独立贡献有限，而 IC50、Se 含量、稳定性与结合能提供了主要增益。